% =====================================================================
%  序：この論文で言うこと，背景，3 つの問いと全体図
%  文の元：paper/main.tex の要旨，paper/sections/1-introduction.tex（貢献 1〜4）
% =====================================================================
\setband{序}

\begin{frame}{要約：少ないラベルは判断基準を先に直し，向きは別のモダリティで補う}
  \rl{要約}
  \vspace{0.1em}
  \begin{itemize}\footnotesize\setlength{\itemsep}{8pt}
        \item \textbf{利得の総当たり実験}：7 群のすべての組で，音声・顔・言葉とその融合を $K$ = 1〜全件で測った（反復 30 回）
        \item \textbf{ラベルが直すもの}：単独の音声と顔の $K\le5$ では，上乗せの \NUM{p1-e29-critshare-kle5-single-lo}〜\NUM{p1-e29-critshare-kle5-single-hi} が\crit{判断基準}の分．音声の\disc{向き}は $q/K$（$q$ = \NUM{p1-qdir-audio-lo}〜\NUM{p1-qdir-audio-hi}）でしか良くならない
        \item \textbf{ソースの価値}：音声の 1 ソースは感情あたり \NUM{p1-e32-n1src-audio-lo}〜\NUM{p1-e32-n1src-audio-hi} 本，全ソースで \NUM{p1-e32-nall-audio-lo}〜\NUM{p1-e32-nall-audio-hi} 本のラベルに当たる
        \item \textbf{手法}：ソースだけで学んだ別のモダリティの分類器（融合の出発点など）がプール全体に弱いラベルを付ける．基本版は 22 の組み合わせのうち \NUM{cmp-ours-vs-best-pos-n} で，拡張版は全部で，最も強い比較手法を上回る
      \end{itemize}
  \takeaway{凍結した特徴＋線形分類器で，理論の予言を実測し，その答えから手法を作った}
\end{frame}

\begin{frame}{新しいコーパスではラベルが少なく，ほかのコーパスを混ぜると得にも損にもなる}
  \rl{背景}
  \vspace{0.2em}
  \begin{columns}[T,onlytextwidth]
    \begin{column}{0.47\textwidth}
      \centering
      \begin{tikzpicture}[font=\scriptsize]
        \node[srcbox, text width=2.3cm] (s1) at (0,1.9) {ソース 1\\（ラベル付き）};
        \node[srcbox, text width=2.3cm] (s2) at (0,0.95) {ソース 2\\（ラベル付き）};
        \node[font=\scriptsize] at (0,0.35) {$\vdots$};
        \node[tgtbox, text width=2.3cm, minimum height=1.25cm] (t) at (0,-0.75) {ターゲット\\感情ごとに $K$ 本の\\ラベル＋ラベル無し};
        \node[clf, text width=1.8cm, minimum height=1.0cm] (c) at (3.2,0.55) {分類器\\（4 感情）};
        \draw[arr] (s1.east) -- (c.west);
        \draw[arr] (s2.east) -- (c.west);
        \draw[arr] (t.east) -- (c.west);
        \node[testbox, text width=1.8cm] (te) at (3.2,-0.75) {ターゲットの\\別の話者で評価};
        \draw[tarr] (c.south) -- (te.north);
      \end{tikzpicture}
    \end{column}
    \begin{column}{0.51\textwidth}
      \begin{itemize}\footnotesize\setlength{\itemsep}{8pt}
        \item 学習に無いコーパスでは，感情の分類器の正解率が下がる
        \item 実務では，ターゲットに感情ごとに数本のラベルを付け，ほかのコーパス（ソース）の発話を足す
        \item 理論は「ラベルが少ないと混ぜて得，多いと損」と言うが，分類器のどこを直すかは言わない
        \item ターゲットの発話には，音声・顔・言葉が同時にあることが多い
      \end{itemize}
    \end{column}
  \end{columns}
\end{frame}

\begin{frame}{3 つの問い：ラベルは何を直すか，ソースは何本ぶんか，その答えで手法を作れるか}
  \rl{問い}
  \vspace{0.4em}
  \begin{itemize}\small\setlength{\itemsep}{10pt}
    \item \textbf{問い 1}　少ない $K$ 本のラベルは，分類器の何を直すか\par
      \muted{\footnotesize → 第 I 部：判断基準と向きに分けて測る}
    \item \textbf{問い 2}　ソースは，ラベルの代わりに何本ぶんを補うか\par
      \muted{\footnotesize → 第 I 部：ソースの向きをラベルの本数に換算する}
    \item \textbf{問い 3}　その答えから，同じ手順で調整した標準の手法を上回る方法を作れるか\par
      \muted{\footnotesize → 第 II 部と比較：残った向きを，別のモダリティの弱いラベルで補う}
  \end{itemize}
  \vspace{0.3em}
  \takeaway{分類器はすべて，凍結したエンコーダの特徴に載せた線形分類器}
\end{frame}

\begin{frame}{全体図：ソースで出発点を作り，$K$ 本は判断基準を直し，弱いラベルが向きを補う}
  \rl{全体図}
  \vspace{-0.3em}
  \centering
  \begin{tikzpicture}[x=0.87cm, y=1.0cm, font=\tiny]
    % 枠
    \fill[cStart!12, rounded corners=3pt] (0,0) rectangle (4.3,5.3);
    \fill[cOurs!8, rounded corners=3pt] (4.5,0) rectangle (10.6,5.3);
    \fill[cExt!8, rounded corners=3pt] (10.8,0) rectangle (17.6,5.3);
    \node[anchor=north west, font=\scriptsize\bfseries] at (0.15,5.15) {ソースのコーパス};
    \node[anchor=north west, text=mMuted] at (0.15,4.72) {ラベル付き．エンコーダは凍結};
    \node[anchor=north west, font=\scriptsize\bfseries] at (4.65,5.15) {第 I 部：$K$ 本が直すもの};
    \node[anchor=north west, text=mMuted] at (4.65,4.72) {単独の A・V，$K\le5$．\textcolor{cOurs}{$\bullet$} ラベルの付いた $K$ 本};
    \node[anchor=north west, font=\scriptsize\bfseries] at (10.95,5.15) {第 II 部：モダリティをまたぐ弱いラベル};
    \node[anchor=north west, text=mMuted] at (10.95,4.72) {テストでは，生徒は自分のモダリティだけを使う};
    % 左：ソースと出発点
    \foreach \y in {3.75,2.9,2.05} {
      \node[draw=cStart, fill=cStart!25, rounded corners=1pt, minimum width=0.75cm, minimum height=0.45cm] at (0.9,\y) {};
      \node[draw=mMuted, fill=white, minimum size=0.32cm, inner sep=0pt, font=\tiny\bfseries] at (1.85,\y) {A};
      \node[draw=mMuted, fill=white, minimum size=0.32cm, inner sep=0pt, font=\tiny\bfseries] at (2.3,\y) {V};
      \node[draw=mMuted, fill=white, minimum size=0.32cm, inner sep=0pt, font=\tiny\bfseries] at (2.75,\y) {T};
    }
    \node[text=mMuted] at (0.9,1.55) {$\vdots$};
    \node[draw=cStart, fill=white, rounded corners=2pt, align=center, text width=2.9cm] (st) at (2.15,0.65) {\textbf{出発点}：全ソースで学んだ\\線形分類器 $(W_0, b_0)$};
    % 中：2 つの感情の雲と境界
    \fill[mMuted!18] (6.3,1.75) ellipse [x radius=0.95cm, y radius=0.55cm, rotate=21];
    \fill[mMuted!18] (8.7,2.75) ellipse [x radius=0.95cm, y radius=0.55cm, rotate=21];
    \node[text=mMuted] at (5.35,2.55) {感情 $k$};
    \node[text=mMuted] at (9.9,2.05) {感情 $l$};
    \pgfmathsetseed{7}
    \foreach \i in {1,...,22} {
      \pgfmathsetmacro{\u}{rand*0.8}\pgfmathsetmacro{\v}{rand*0.38}
      \fill[mMuted!75] ($(6.3,1.75)+(21:\u)+(111:\v)$) circle (0.7pt);
      \pgfmathsetmacro{\u}{rand*0.8}\pgfmathsetmacro{\v}{rand*0.38}
      \fill[mMuted!75] ($(8.7,2.75)+(21:\u)+(111:\v)$) circle (0.7pt);
    }
    \foreach \p in {(5.95,1.85),(6.6,1.6),(6.25,2.1)} \fill[cOurs] \p circle (1.6pt);
    \foreach \p in {(8.35,2.9),(9.05,2.6),(8.75,3.05)} \fill[cOurs] \p circle (1.6pt);
    \draw[cStart, line width=1.0pt, dash pattern=on 3pt off 1.5pt] (7.25,0.5) -- (6.747,3.35);
    \draw[cOurs, line width=1.0pt] (7.809,0.5) -- (7.306,3.35);
    \draw[mDarkTeal, line width=0.7pt, dotted] (8.172,0.5) -- (7.078,3.35);
    \draw[->, cCrit, line width=0.8pt] (7.18,0.9) -- (7.74,0.9);
    \node[anchor=west, text=cCrit, align=left] at (8.3,0.75) {判断基準 $o$\\誤差 $1/(2K)$\\先に直る};
    \draw[->, cDisc, line width=0.6pt] (7.5,2.25) ++(100:1.0) arc[start angle=100, end angle=112, radius=1.0cm];
    \node[anchor=west, text=cDisc, align=left] at (7.6,4.0) {向き（識別力 $d$ を決める）\\誤差 $q/K$，直りにくい};
    \draw[cStart, line width=1.0pt, dash pattern=on 3pt off 1.5pt] (4.7,4.2) -- ++(0.45,0); \node[anchor=west] at (5.15,4.2) {出発点};
    \draw[cOurs, line width=1.0pt] (4.7,3.92) -- ++(0.45,0); \node[anchor=west] at (5.15,3.92) {$K$ 本の後};
    \draw[mDarkTeal, line width=0.7pt, dotted] (4.7,3.64) -- ++(0.45,0); \node[anchor=west] at (5.15,3.64) {ターゲットで最良};
    \draw[arr, cStart] (st.east) -- (5.6,0.65);
    % 右：先生，弱いラベル，生徒
    \node[draw=cExt, fill=white, rounded corners=2pt, align=center, text width=2.1cm] (tea) at (12.2,3.1) {\textbf{先生}：融合の\\分類器（A+V+T）\\ソースだけで学ぶ};
    \begin{scope}[shift={(13.95,3.55)}]
      \foreach \r/\a/\b/\c/\d in {0/60/15/10/15, 1/10/55/20/15, 2/15/10/65/10, 3/20/20/10/50, 4/45/25/15/15} {
        \fill[cExt!\a] (0.00,-\r*0.27) rectangle ++(0.36,0.24);
        \fill[cExt!\b] (0.39,-\r*0.27) rectangle ++(0.36,0.24);
        \fill[cExt!\c] (0.78,-\r*0.27) rectangle ++(0.36,0.24);
        \fill[cExt!\d] (1.17,-\r*0.27) rectangle ++(0.36,0.24);
      }
      \node[anchor=south] at (0.77,0.24) {4 感情};
      \draw[decorate, decoration={brace, amplitude=2.5pt}, mMuted] (1.62,0.24) -- (1.62,-1.08);
      \node[anchor=west, align=left] at (1.72,-0.42) {ラベル無しの\\発話すべて};
    \end{scope}
    \draw[arr, cExt] (tea.east) -- (13.9,3.1);
    \node[anchor=west, text=cExt, align=left] at (13.9,2.05) {弱いラベル：感情ごとの\\合計をそろえる};
    \draw[arr, cExt] (15.6,1.75) -- (16.1,1.0);
    \node[draw=cOurs, fill=white, rounded corners=2pt, align=center, text width=1.8cm] (stu) at (16.35,0.6) {\textbf{生徒}：線形\\1 つのモダリティ};
    \node[draw=cStart, fill=white, rounded corners=2pt, align=center, text width=1.6cm] (st2) at (12.0,0.6) {出発点\\（同じモダリティ）};
    \draw[cStart, line width=0.9pt, decorate, decoration={coil, aspect=0.4, segment length=3pt, amplitude=2.2pt}] (st2.east) -- (stu.west);
    \node[anchor=south] at ($(st2.east)!0.5!(stu.west)+(0,0.12)$) {出発点に縛る};
    \foreach \x in {0,1,2} \fill[cOurs] ($(11.25,1.85)+(\x*0.2,0)$) circle (1.6pt);
    \node[anchor=west, text=cOurs] at (11.8,1.85) {$K$ 本};
    \draw[arr, cOurs] (11.45,1.65) .. controls (11.45,1.35) and (13.0,1.45) .. (14.4,1.15);
  \end{tikzpicture}\par
  \vspace{0.2em}
  {\raggedright\scriptsize 中の図は模式図．$K$ 本の後（青）は出発点（灰）を平行に動かすが，最良の向き（点線）までは回らない\par}
\end{frame}
